\section{Comparison with Other Formalizations}\label{sec:related}

We compare along the three axes the library combines --- (i) presentations
and normal forms of groups and monoids, (ii) completeness of equational
theories, (iii) circuits --- across Agda, Lean, Coq, and Isabelle.
Table~\ref{tab:compare} summarizes; the prose gives the details and the
precise deltas.  Claims of absence are stated to the best of our
knowledge.

\begin{table}
\caption{Mechanized work adjacent to this library.  ``Pres.''\ = concrete
presentations by generators and relations; ``NF'' = verified normal forms
for presented structures; ``Compl.''\ = completeness of a relation set for
an independent semantics.}
\label{tab:compare}
\centering
\small
\begin{tabular}{@{}llccccl@{}}
\toprule
Development & Prover & Circuits & Pres. & NF & Compl. & Focus \\
\midrule
This library & Agda & \ding{51} & \ding{51} & \ding{51} & \ding{51} & gate groups, constructions \\
Bian--Selinger devs.~\cite{bian2023cliffordt2,bian2023cliffordcs3} & Agda & \ding{51} & \ding{51} & \ding{51} & syntactic & two presentations relate \\
\textsc{sqir}/\textsc{voqc}~\cite{hietala2021voqc} & Coq & \ding{51} & -- & -- & sound only & optimizer correctness \\
\textsc{Qwire}~\cite{paykin2017qwire} & Coq & \ding{51} & -- & -- & -- & circuit language \\
VyZX~\cite{lehmann2025vyzx} & Coq & \ding{51} & -- & -- & sound only & ZX diagram rules \\
CoqQ~\cite{zhou2023coqq} & Coq & \ding{51} & -- & -- & -- & program logics \\
Qbricks~\cite{chareton2021qbricks} & Why3 & \ding{51} & -- & -- & -- & automated deduction \\
QHLProver~\cite{liu2019quantum}, IMD~\cite{bordg2020dirac} & Isabelle & \ding{51} & -- & -- & -- & algorithms, logics \\
mathlib groups~\cite{mathlib2020,mathlib-doc-presented} & Lean & -- & \ding{51} & -- & -- & abstract theory \\
Odd Order~\cite{gonthier2013oddorder} & Coq & -- & -- & -- & -- & finite group theory \\
Free groups~\cite{breitner2010freegroups,kharim2023freegroups} & Isabelle & -- & \ding{51} & \ding{51} & -- & Nielsen--Schreier \\
IsaFoR completion~\cite{sternagel2013kb,hirokawa2019abstract} & Isabelle & -- & -- & \ding{51} & -- & term rewriting \\
cubical/unimath~\cite{cubicallibrary,agdaunimath,mangel2023sign} & Agda & -- & partial & -- & -- & univalent groups \\
\bottomrule
\end{tabular}
\end{table}

\subsection{The Bian--Selinger verifications: syntax to syntax}

The direct ancestors of this library are the Agda developments that Bian
and Selinger published alongside their presentations of the $2$-qubit
Clifford+$T$ group~\cite{bian2023cliffordt2} and the $3$-qubit
Clifford+CS group~\cite{bian2023cliffordcs3}, systematized in Bian's
thesis~\cite{bian2023thesis}.  Methodologically these pioneered
everything we build on: monoid-level presentations (no formal inverses),
the Reidemeister--Schreier theorem applied --- to Greylyn's presentation
of $U_4(\mathbb{Z}[1/\sqrt2,i])$~\cite{greylyn2014} in the $2$-qubit
case, recursively in the CS case --- to derive subgroup presentations,
amalgamated products of monoids, and machine-checked simplification of
``thousands of relations'' down to a small set (seventeen, for
Clifford+CS).  The crucial
limitation, called out in our introduction, is that these developments
work entirely on the \emph{syntactic side}: the theorems relate two
\emph{presentations} of a group --- the isomorphism of two finitely
presented monoids --- and no independent semantic object (a matrix group,
a permutation group, a product of groups) occurs in any statement.
Soundness and completeness in the model-theoretic sense are therefore not
expressible there, only relative completeness of one relation set against
another.  The present library keeps the syntactic technology (our
amalgamation case studies end in exactly such isomorphisms) and adds the
semantic half: the \aid{{\_}IsPresentationOf{\_}} records of
Table~\ref{tab:presentations} target groups built independently of the
syntax --- stdlib permutations, \texttt{ForStdlib} products and
extensions --- and the generic \aid{by-normalization} lemma is what
connects the two halves.  A second delta is reusability: coset tables,
towers, packed levels, and the product constructions are library modules
with explicit hypothesis records, where the earlier developments were
structured as per-paper proofs.

\subsection{Quantum circuits in proof assistants: soundness, not completeness}

A rich ecosystem verifies quantum \emph{programs} and \emph{compilers}.
\textsc{Qwire}~\cite{paykin2017qwire} embeds a circuit language in Coq
with a denotational semantics; \textsc{sqir}/\textsc{voqc}
\cite{hietala2021voqc,hietala2021proving} defines a quantum IR in Coq and
verifies an optimizing compiler --- every rewrite the optimizer performs
is proved \emph{sound} against the matrix semantics.
VyZX~\cite{lehmann2025vyzx} inductively represents ZX diagrams in Coq and
proves the ZX rewrite rules sound.  CoqQ~\cite{zhou2023coqq} builds
Dirac-notation program logics over MathComp;
Qbricks~\cite{chareton2021qbricks} verifies circuit-building programs
with automated first-order reasoning over path-sums; on the Isabelle
side, QHLProver formalizes quantum Hoare logic~\cite{liu2019quantum} and
Isabelle-Marries-Dirac verifies textbook algorithms~\cite{bordg2020dirac}.
None of these address the question this paper is about: whether a
\emph{finite rule set derives all true equations} --- the completeness of
the calculus itself, as opposed to the correctness of particular
transformations in it.  The distinction is not academic: completeness is
what justifies normal-form-based decision procedures and tells an
optimizer author that no valid rewrite is missing.  The mathematical
literature proves such theorems on
paper~\cite{selinger2015clifford,clement2023complete,clement2024minimal,
li2025qutrit,clement2026realcliffordch} --- likewise for the graphical
side, where the completeness results for the ZX-calculus, from the qubit
stabilizer fragment through Clifford+$T$ to the odd-prime-dimensional
qupit calculi~\cite{backens2014zx,jeandel2018zx,booth2022qupitzx,
poor2023travaganza}, are all pen-and-paper; mechanizations have so far
been confined to the Bian--Selinger line and the present work.  We see the two
ecosystems as complementary: a natural target is to connect our presented
gate groups to the matrix semantics of \textsc{sqir} or VyZX, obtaining
end-to-end completeness against matrices rather than against abstract
product groups.

\subsection{Group theory in proof assistants: abstract theory versus concrete presentations}

Lean's mathlib~\cite{mathlib2020} defines \aid{FreeGroup} and
\aid{PresentedGroup} --- the quotient of a free group by the normal
closure of a relator set --- and develops Coxeter groups atop
them~\cite{mathlib-doc-presented}, together with a groupoid-based proof
of the Nielsen--Schreier theorem.  This is presentation \emph{theory} in
the abstract: what mathlib does not provide is infrastructure for
proving that a \emph{given} finite relation set presents a \emph{given}
concrete group --- no coset enumeration, no verified Schreier
transversals, no normal-form transports; the Coxeter development, for
instance, axiomatizes the presentation rather than verifying completeness
of specific relation sets for specific permutation groups.  Conversely,
the deep results of the Coq/MathComp school --- most famously the Odd
Order theorem~\cite{gonthier2013oddorder} --- concern finite groups
represented concretely (as subgroups of permutation groups over finite
types), and presentations play no role.  In Isabelle, the AFP has free
groups with the ping-pong lemma~\cite{breitner2010freegroups} and a
recent combinatorial Nielsen--Schreier proof with a decision procedure
for the conjugacy problem in free groups~\cite{kharim2023freegroups} ---
normal-form reasoning in the free group, adjacent in spirit to our
\aid{NormalFormInjective}, but again without presentations of concrete
target groups.  Against this landscape the library's niche is sharp:
\emph{verified coset enumeration for concrete presentations}, of the kind
computational group theory performs unverified~\cite{sims1994computation,
holt2005handbook}, packaged so that the certificate (tables, sections,
five equations) is separated from the search that found it.

On the univalent side, cubical Agda~\cite{vezzosi2019cubical} and its
library~\cite{cubicallibrary} provide free groups as higher inductive
types and quotients with the right universal properties definitionally;
agda-unimath develops finite group theory including the symmetric groups
and the sign homomorphism~\cite{agdaunimath,mangel2023sign}.  A HIT-based
presented group avoids the congruence bookkeeping our setoid design pays
--- functions out of the quotient are just functions --- but our central
devices are \emph{maps into} the syntax (sections of normal forms,
Schreier transversals), which quotients complicate: a section
$\mathit{NF} \to \aid{Word}\;X$ picking canonical representatives is
setoid-natural and quotient-hostile.  Since the library also targets the
(non-cubical) standard library and must remain
\texttt{--cubical-compatible} rather than cubical, the setoid design
is the pragmatic point in the space; \S\ref{sec:discussion} reflects
further.

\subsection{Formalized rewriting and completion}

The rewriting community has formalized its own meta-theory: IsaFoR/CeTA
contains Knuth--Bendix orders and completion~\cite{sternagel2013kb} and
abstract completion~\cite{hirokawa2019abstract}, certifying completion
\emph{procedures} that, when they succeed, decide word problems via
convergent rewriting.  Our approach is deliberately different: we never
orient the relations or prove confluence --- the normal form is an
arbitrary verified \emph{function}, and uniqueness is proved against the
semantics.  This sidesteps the classical obstruction that finitely
presented monoids with decidable word problem need not admit finite
convergent presentations (Squier~\cite{squier1987word,
squier1994finiteness}) --- a normal-form function is a weaker artifact
than a convergent rewrite system, and correspondingly more often
available.  The polygraphic school of Guiraud--Malbos and
collaborators~\cite{guiraud2018polygraphs,ara2025polygraphs} studies
presentations and their coherence in higher categories, with Lafont's
algebraic theory of Boolean circuits~\cite{lafont2003boolean} as an early
landmark directly ancestral to circuit presentations like ours; those
developments are pen-and-paper, and mechanizing polygraphic coherence
along the lines of this library is an attractive direction.

\subsection{Summary of the niche}

To our knowledge, this library provides: the first mechanized
\emph{semantic} completeness proofs for circuit relation sets (symmetric
groups for two semantics; wreath and Pauli families; with the qutrit
Clifford+$T$ presentation, the first machine-checked presentation
theorem of any qutrit gate monoid); the first verified
Reidemeister--Schreier/coset-enumeration engine and amalgamated-product
normal forms in any proof assistant; and the first presentation-level
product/extension theorems connected to standard-library-style semantic
constructions.  Each claim is scoped by Table~\ref{tab:compare} and the
paragraphs above.
